\chapter{Методика}
\label{ch:methods}

%надо провести адиабатический процесс, чтобы была связь с помощью уравнения Пуассона. Есть сосуд фиксированного объёма на основе этого создать адиабатический процесс, на основе этого описать методику
%входное выходное давление. как сделать так, чтобы температуру не учитывать, вводим изотермический процесс
Пусть имеется сосуд фиксированного размера. Для определения отношения $C_p/C_v$ необходимо провести адиабатический процесс, подчиняющийся уравнению Пуассона.

Создав избыточное давление, можно повысить температуру газа по отношению к окружающей среде. Вследствие теплообмена со стенками сосуда начнётся изохорное охлаждение. Через некоторое время газ остынет до комнатной температуры. При этом давление воздуха понизится до $p_0 + \Delta p_1$, где

\begin{equation}
    \Delta p_1 = \rho g \Delta h_1.
\end{equation}

Открыв кран, соединяющий сосуд с окружающей средой, можно запустить адиабатическое расширение газа, вследствие чего его температура окажется ниже комнатной. Через время порядка $0.5$ с газ начнёт изобарически нагреваться. Выберем время $\tau$, в течение которого кран остаётся открытым, таким, чтобы можно было пренебречь временем адиабатического расширения воздуха. После закрытия крана газ станет изохорически нагреваться до комнатной температуры, причём давление внутри сосуда возрастёт до $p_0 + \Delta p_2,$ где

\begin{equation}
    \Delta p_2 = \rho g \Delta h_2.
\end{equation}

Для описания состояния газа можно использовать модель идеального газа. Тогда рассмотрев изобарическое расширение, используя уравнение теплового баланса для имеющейся во времени массы газа $m = \frac{p_0V_0}{RT}\mu$:

\begin{equation}
    c_p m \ dT = -\alpha (T - T_0) dt, 
\end{equation}

где $c_v$ — удельная теплоёмкость воздуха при постоянном давлении, $a$ — положительный постоянный коэффициент, характеризующий теплообмен, $V$ — объем сосуда, после преобразований с использованием выражений (1) и (2) можно получить выражение для определения показателя адиабаты $\gamma$ (Вывод см. в Приложении А.):

\begin{equation}
    \ln{\frac{\Delta h_1}{\Delta h_2}} = \ln{\frac{\gamma}{\gamma - 1}} + \left( \frac{\alpha}{c_p m_0}  \tau\right) 
\end{equation}

Тогда исследуя зависимость отношения перепадов давления $\frac{\Delta p_1}{\Delta p_2}$ от времени $\tau$, можно определить показатель адиабаты. Например, его можно найти из графика зависимости $\ln{\frac{\Delta h_1}{\Delta h_2}}$ от $\tau$.

